
\documentclass[conference]{IEEEtran}
% \IEEEoverridecommandlockouts
% The preceding line is only needed to identify funding in the first footnote. If that is unneeded, please comment it out.
%Template version as of 6/27/2024

\usepackage{cite}
\usepackage{amsmath}
\usepackage{amssymb}
\usepackage{amsfonts}
\usepackage{algorithmic}
\usepackage{graphicx}
\usepackage{textcomp}
\usepackage{xcolor}
\usepackage{newtxtext}
\usepackage{enumitem}
\usepackage{url}

\def\BibTeX{{\rm B\kern-.05em{\sc i\kern-.025em b}\kern-.08em
    T\kern-.1667em\lower.7ex\hbox{E}\kern-.125emX}}
\begin{document}

\title{Seminar Paper on APL Ethernet}
\author{
    \IEEEauthorblockN{1\textsuperscript{st} Raihan Uddin}
    \IEEEauthorblockA{
        \textit{Electronic Engineering} \\
        \textit{Student ID: 2210018}\\
        \textit{Hochschule Hamm-Lippstadt}\\
        Lippstadt, Germany \\
        raihan.uddin@stud.hshl.de
    }
    \and
\IEEEauthorblockN{2\textsuperscript{nd} Bodrul Amin Murad}
\IEEEauthorblockA{\textit{Electronic Engineering} \\
\textit{Student ID: 2210039}\\
\textit{Hochschule Hamm-Lippstadt}\\
Lippstadt, Germany \\
bodrul-amin.murad@stud.hshl.de}
\and
\IEEEauthorblockN{3\textsuperscript{rd} Ismail Hossain}
\IEEEauthorblockA{\textit{Electronic Engineering} \\
\textit{Student ID: 2210015}\\
\textit{Hochschule Hamm-Lippstadt}\\
Lippstadt, Germany \\
ismail.hossain@stud.hshl.de}
}
\maketitle

\begin{abstract}
%APL is a new physical layer of networking, based on existing Ethernet standard. The increasing use of Industrial Ethernet can add more values and efficiency if they can connect all devices and can be replaced as a single homogeneous solution in the field, instead of multiple heterogeneous protocols. 

\end{abstract}

\begin{IEEEkeywords}
%Ethernet, networking, standards, cable.
\end{IEEEkeywords}
\section{Introduction}
The Cross Traffic Management project deals with the complexity of efficient routing of autonomous vehicles through a four-way intersection. Conventional traffic control methods lack the adaptability and responsiveness needed to manage real-time vehicle flows for autonomous vehicles. To avert these limitations, our implemented system, here, employs a max dynamic weight scheduling algorithm--- dynamically assigning priority to each traffic direction based on real-time queue lengths, vehicle arrival rates, and safety constraints.
\begin{enumerate}[label=(\alph*)]
    \item \textbf{Core Logic in C:} The foundational behavior of the intersection controller was implemented in C. This stage established the data structures and algorithms needed to model vehicle queues and calculate dynamic weights for each direction.
    \item \textbf{FreeRTOS Integration:} To support real-time execution and concurrent processing of sensor data, queue management, and scheduling decisions, the C implementation was ported to FreeRTOS. This ensured deterministic task scheduling and inter-task communication under realistic load conditions.
    \item \textbf{UPPAAL Verification:} Using UPPAAL, we formally modeled the system’s timing behavior and verified critical properties such as deadlock freedom, bounded waiting time, and correct weight recalculation. This step provided mathematical guarantees of system correctness before hardware realization.
    \item \textbf{VHDL Hardware Design:} Finally, the scheduling logic was translated into VHDL and simulated in ModelSim. This hardware-oriented implementation demonstrates the feasibility of deploying the controller on FPGA or ASIC platforms for ultra-low-latency intersection management.
\end{enumerate}
To achieve the objectives of efficient cross-traffic management, our work is structured around the following research questions:
\begin{enumerate}[label=(\alph*)]
    \item \textbf{Queuing Model:} How can we accurately model intersection traffic for autonomous vehicles using queuing theory?
    \item \textbf{Optimal Scenario:} What scheduling scenarios and weight functions best minimize congestion and waiting times?
    \item \textbf{Simulation Framework:} How can we simulate the intersection controller—first in C and then under FreeRTOS—and which parameters (e.g., arrival rates, service times, task priorities) are critical to assess performance?
    \item \textbf{Hardware Realization:} Which components of the system are best suited for hardware implementation, and how can VHDL simulation validate their real-time behavior?
\end{enumerate}
By combining software development, real-time operating systems, formal verification, and digital circuit design, this project delivers a scalable and robust traffic control solution suitable for smart intersections in autonomous vehicle networks.
\section{Scheduling Algorithm}
In our centralized traffic management system, we have used dynamic-weight scheduling algorithm to efficiently route autonomous vehicles at the junction. For algorithm design, we sought online help that is browsing for critical steps in ChatGpt. According to suggested guidelines, we designed our algorithm as follows:
\subsection{Overview and Objectives}
The algorithm is meant to achieve:
\begin{enumerate}[label=(\alph*)]
    \item Dynamic prioritization of competing traffic directions based on real-time metrics.
    \item Fairness (no direction is starved) and throughput optimization (minimize overall waiting time).
    \item Safety constraints, e.g., minimum green times, no conflicting movements.
\end{enumerate}
\subsection{System Model and Metrics}
The input and state variables of the scheduler are:
\begin{enumerate}[label=(\alph*)]
    \item Queue length $Q_i$ for each direction $i \in \{N, S, E, W\}$
    \item Arrival rate $\lambda_i$
 (vehicles/sec) estimated over a sliding window
    \item Waiting time $W_i$ of the lead vehicle in the queue
    \item Safety timer $T_{\text{min}}$ ensuring a green phase lasts at least some minimum
\end{enumerate}
\subsection{Weight Computation}
Present the weight function that combines these metrics into a single score per direction:
$w_i$= $\alpha$$Q_i$ +  $\beta$$\lambda_i$ + $\gamma$$W_i$\\
where, $\alpha$, $\beta$ and $\gamma$ are tunable coefficients that reflect the relative importance of queue length, flow rate and waiting time.
\subsection{Algorithm Steps}
The loop of the scheduler will work like follows:
\begin{enumerate}
    \item Sense and Update
        \begin{enumerate}[label=(\alph*)]
            \item Read current queue lengths $Q_i$ and arrival rates of $\lambda_i$.
            \item Update waiting times $W_i$ for vehicles in the queue.
        \end{enumerate}
    \item Compute weights
        \begin{enumerate}[label=(\alph*)]
            \item For each direction i, we need to compute $w_i$.
        \end{enumerate}  
    \item Select Max-Weight Direction
            \begin{enumerate}[label=(\alph*)]
            \item We have to select \[d = \arg\max_i (w_i)\]
            \item Enforce safety constraints, if the previously green directions has not yet met $T_{\text{min}}$, delay switching.
            \end{enumerate}  
    \item Grant Green phase
            \begin{enumerate}[label=(\alph*)]
            \item Activate green for direction d.
            \item Start or reset the green-phase timer $T_d$
            \end{enumerate}  
    \item Repeat
            \begin{enumerate}[label=(\alph*)]
                \item After either a fixed epoch $T_{\text{max}}$ or when $w_d$ falls below the next-highest weight by a threshold, return to the initial point.
            \end{enumerate}  
    \item Properties and Tuning
            \begin{enumerate}[label=(\alph*)]
                \item \textbf{Fairness:}Because every direction’s weight eventually grows (via $W_i$), starvation is prevented.
                \item \textbf{Responsiveness:}By tuning $\alpha$, $\beta$ and $\gamma$, the system can emphasize clearing long queues or reacting to sudden surges.
                \item \textbf{Safety Margins:}$T_{\text{min}}$ avoids rapid toggling.
            \end{enumerate}
\end{enumerate}
\begin{figure}[htbp]
    \centering
    \includegraphics[width=0.9\columnwidth]{C:/Users/exone/Documents/GitHub/Embedded_Engineering_Lab/Sections/diagrams/Q0.jpeg}
    \caption{Q0}
    \label{fig:Q0}
\end{figure} 
\begin{figure}[htbp]
    \centering
    \includegraphics[width=0.9\columnwidth]{C:/Users/exone/Documents/GitHub/Embedded_Engineering_Lab/Sections/diagrams/Q1.jpeg}
    \caption{Q1}
    \label{fig:Q1}
\end{figure} 
\begin{figure}[htbp]
    \centering
    \includegraphics[width=0.9\columnwidth]{C:/Users/exone/Documents/GitHub/Embedded_Engineering_Lab/Sections/diagrams/Q2.jpeg}
    \caption{Q2}
    \label{fig:Q2}
\end{figure} 
\begin{figure}[htbp]
    \centering
    \includegraphics[width=0.9\columnwidth]{C:/Users/exone/Documents/GitHub/Embedded_Engineering_Lab/Sections/diagrams/Q3.jpeg}
    \caption{Q3}
    \label{fig:Q3}
\end{figure} 
\section{SysML Diagram}
We have used different SysML diagrams to visually represent the system's structure and behavior.
\subsection{Activity Diagram}
This diagram models the runtime behavior of a single vehicle as it interacts with the intersection controller:
\begin{figure}[htbp]
    \centering
    \includegraphics[width=0.9\columnwidth]{C:/Users/exone/Documents/GitHub/Embedded_Engineering_Lab/Diagrams/Activity_diagram.png}
    \caption{Activity diagram}
    \label{fig:Activity diagram}
\end{figure} 
\begin{enumerate}
    \item Receive arrival + turn intent
        \begin{enumerate}[label=(\alph*)]
            \item A vehicle arrives at the stop line and communicates its intended movement (straight, left, or right).
        \end{enumerate}
    \item Determine bundle
        \begin{enumerate}[label=(\alph*)]
            \item Based on direction (North South or East West) and turn intent (through or left), the vehicle is assigned to one of four “bundles”: NS-TR, EW-TR, NS-L, or EW-L.
        \end{enumerate}  
    \item Enqueue into bundle queue
            \begin{enumerate}[label=(\alph*)]
            \item The controller places the request in the corresponding queue for later scheduling.
            \end{enumerate}  
    \item Check against current green phase
            \begin{enumerate}[label=(\alph*)]
                \item If the vehicle bundle matches the currently active (green) phase, it bypasses waiting; otherwise:
                \begin{enumerate}[label=(\alph*)]
                \item Send “WAIT” signal
                \item Remain queued until the green phase switches back.
                \end{enumerate}  
            \end{enumerate}  
    \item Send “GO” signal
            \begin{enumerate}[label=(\alph*)]
                \item When its bundle is selected by the scheduling algorithm, the controller grants green.
            \end{enumerate}  
    \item Cross intersection safely
            \begin{enumerate}[label=(\alph*)]
                \item The vehicle proceeds through the junction under the granted green signal.
            \end{enumerate}
\end{enumerate}

\subsection{Requirement Diagram}
This UML requirements diagram organizes and decomposes the system-level requirement “Intersection Control” into five constituent requirements:
\begin{figure}[htbp]
    \centering
    \includegraphics[width=0.9\columnwidth]{C:/Users/exone/Documents/GitHub/Embedded_Engineering_Lab/Diagrams/requirement_diagram.jpg}
    \caption{Requirement diagram}
    \label{fig:Requirement diagram}
\end{figure} 
\begin{enumerate}
    \item Vehicle Detection
        \begin{enumerate}[label=(\alph*)]
            \item``The system shall detect and count vehicles approaching from all directions''.
        \end{enumerate}
    \item Movement Schedule
        \begin{enumerate}[label=(\alph*)]
            \item``The system shall schedule non-conflicting movements to maximize parallel throughput''.
        \end{enumerate}  
    \item Central Hub Communication
            \begin{enumerate}[label=(\alph*)]
            \item ``The system shall send queue information to and receive grants from a central controller''.
            \end{enumerate}  
    \item Parameter Configuration
            \begin{enumerate}[label=(\alph*)]
                \item ``The system shall allow configuration of timing parameters (e.g., green/gap duration)''.
            \end{enumerate}  
    \item Safety Enforcement
            \begin{enumerate}[label=(\alph*)]
                \item ``The system shall enforce all-red gaps and prevent any conflicting movements''.
            \end{enumerate}
\end{enumerate}

\subsection{Use Case Diagram}
This diagram captures the key interactions between two actors and the Infrastructure Intersection Management System:
\begin{figure}[htbp]
    \centering
    \includegraphics[width=0.9\columnwidth]{C:/Users/exone/Documents/GitHub/Embedded_Engineering_Lab/Diagrams/use_case.png}
    \caption{Requirement diagram}
    \label{fig:Requirement diagram}
\end{figure} 

\begin{enumerate}
    \item \textbf{Actor:} Car
        \begin{enumerate}[label=(\alph*)]
            \item Enter Queue
            \item Enter/Request to Enter Junction
                \begin{enumerate}[label=(\alph*)]
                    \item Includes the ``Manage Traffic Flow using Queuing Algorith'' use case.
                \end{enumerate}
            \item Exit Junction
        \end{enumerate}  
    \item \textbf{Actor:} Intersection Control
            \begin{enumerate}[label=(\alph*)]
            \item Invokes “Manage Traffic Flow using Queuing Algorithm” (representing the central controller or coordinating entity that runs the scheduling logic).
            \end{enumerate}  
\end{enumerate}


\newpage
\bibliographystyle{IEEEtran}  % Change to apa, plain, acm, etc. if required
%\bibliography{bibliography_Comm_standard}  % Your .bib file (no .bib extension)

\end{document}
